\documentclass[11pt]{article}
\usepackage[margin=1in]{geometry}
\usepackage{amsmath, amssymb, amsthm, mathtools}
\usepackage{stmaryrd}
\usepackage{bussproofs}
\usepackage{hyperref}
\usepackage{listings}
\usepackage{xcolor}
\usepackage{enumitem}

\definecolor{rhocolor}{HTML}{1A1A2E}
\definecolor{rhoaccent}{HTML}{3FA9F5}

\lstdefinelanguage{rholang}{
  morekeywords={new, in, for, contract, match, if, else, true, false, Nil, send, receive},
  sensitive=true,
  morecomment=[l]{//},
  morestring=[b]"
}
\lstset{
  basicstyle=\ttfamily\small,
  keywordstyle=\color{rhocolor}\bfseries,
  commentstyle=\color{gray}\itshape,
  stringstyle=\color{rhoaccent},
  numbers=left,
  numberstyle=\tiny\color{gray},
  frame=single,
  breaklines=true,
  showstringspaces=false,
  columns=fullflexible
}

\theoremstyle{definition}
\newtheorem{definition}{Definition}[section]
\newtheorem{example}[definition]{Example}
\theoremstyle{plain}
\newtheorem{proposition}[definition]{Proposition}
\newtheorem{theorem}[definition]{Theorem}
\theoremstyle{remark}
\newtheorem*{remark}{Remark}

\newcommand{\Sig}{\mathsf{Sig}}
\newcommand{\Alice}{\mathsf{Alice}}
\newcommand{\Bob}{\mathsf{Bob}}
\newcommand{\Dave}{\mathsf{Dave}}
\newcommand{\sem}[1]{\llbracket #1 \rrbracket}
\newcommand{\Att}{\mathrel{\Vdash}}
\newcommand{\AttD}{\Att_{\!\mathsf{D}}}
\newcommand{\AttW}{\Att_{\!\mathsf{W}}}
\newcommand{\AttH}{\Att_{\!\mathsf{H}}}
\newcommand{\modD}[1]{\langle #1 \rangle^{\mathsf{D}}}
\newcommand{\modW}[1]{\langle #1 \rangle^{\mathsf{W}}}
\newcommand{\modH}[1]{\langle #1 \rangle^{\!\mathsf{H}}}
\newcommand{\Reps}{\mathcal{R}}
\newcommand{\Traces}{\mathcal{T}}
\newcommand{\Proofs}{\mathcal{P}}
\newcommand{\Names}{\mathcal{N}}

\title{Identity-Indexed Typing Judgments\\
and the Adjudication of Capability\\[0.5em]
\large A Proof-Theoretic / Observational Framework\\for Less Game-able Reputation}
\author{L.G.~Meredith \\ \small F1R3FLY Industries \\[0.25em] \small in dialogue with Claude (Anthropic)}
\date{\today}

\begin{document}
\maketitle

\begin{abstract}
We describe a reputation framework in which the term position of a typing judgment is occupied not by inspectable code but by an opaque cryptographic identity. Conformance of an identity to a spatial-behavioral type is adjudicated by the community through two simultaneously running processes: a distributed proof system over signed judgments, and a labelled transition system assembled from attested communication events. Reputation is the lattice position of an identity between proof-theoretic certification and observational refutation, refusing scalarization until a community-specific decision compels projection. The attestation discipline is stratified into direct, witnessed, and hearsay tiers in a way that aligns precisely with the name-as-capability principle of the rho calculus, inducing an evidence-indexed family of Hennessy--Milner modalities through the OSLF construction. We sketch a hosted implementation in MeTTaIL (rholang~1.4) and draw structural parallels with zero-knowledge proof systems, where the witness is likewise hidden from the verifier.
\end{abstract}

\section{Introduction}

\subsection{The inversion}

In conventional type theory, a judgment
\[
\Gamma \vdash M : T
\]
exhibits the term $M$ syntactically. The typing relation is defined by induction over the structure of $M$; type checking presupposes access to the code. The identity of any agent producing or owning $M$ is irrelevant to the judgment. The whole edifice presumes a transparent term-bearer.

In the framework we describe, the term position is occupied instead by a cryptographic identity:
\[
\Gamma \vdash \Sig(\Alice) : T.
\]
The identity is \emph{opaque} --- we have no access to Alice's internal code, no way to perform induction over her behaviour --- but \emph{rigidly locatable}: under standard cryptographic assumptions, signatures verifiably tie messages to identities, and impersonation is bounded by hardness.

This inversion is the load-bearing conceptual maneuver. We are no longer asking whether a given program inhabits a given type. We are asking: given that an identity \emph{claims} to inhabit a type, by what process can the community adjudicate the claim?

A short anticipatory remark. The identity $\Sig(\Alice)$ should not be read as a static, eternal object. It is a \emph{binding} maintained by a persistent system that issues cryptographic identifiers to entering agents, sustains the binding while they participate, and retires or recycles the identifier when they depart. The community is dynamic; what is fixed is this identity registry. We develop the registry-relative reading in \S\ref{ssec:registry}; the rest of the framework rests on it.

\subsection{What spatial-behavioral types contribute}

A spatial-behavioral type, in the sense relevant to the rho calculus and its successors, describes both \emph{what} an agent communicates (the behavioral dimension, in the tradition of session types and process calculi) and \emph{where} the communication is structurally located (the spatial dimension, capturing nominal scope, locality, and the topology of channel sharing). The graph-structured lambda theories (GSLTs) generalise this further: a type is a configuration in a grammar quotiented by equations and oriented by rewrites, and the OSLF functor associates to each such theory a Hennessy--Milner-style logic whose characteristic formulae enjoy an adequacy theorem with respect to bisimulation.

This is what makes spatial-behavioral types suitable as the right-hand side of the identity judgment. They are rich enough to describe \emph{services} --- patterns of interaction the community might wish to reason about --- while admitting a model-theoretic semantics that interacts cleanly with observation. A trace of communication events either falsifies a type by satisfying a refuting formula, or fails to falsify it.

\subsection{Plan}

Section~\ref{sec:framework} develops the dual structure of distributed proof and communal trace, and defines reputation as a lattice position between them. Section~\ref{sec:attest} introduces the three-tier attestation discipline and shows how it aligns with rho's name-as-capability principle. Section~\ref{sec:diff} contrasts the framework with the family of typeless-scalar reputation systems prevalent in practice. Section~\ref{sec:mettail} sketches a hosted implementation in MeTTaIL. Section~\ref{sec:zk} examines the structural parallels with zero-knowledge proof systems and identifies several places where ZK techniques compose naturally with the framework.

\section{The framework}\label{sec:framework}

\subsection{Identity-indexed judgments}

We assume an identity registry $\mathcal{R}$ that issues cryptographic identifiers to agents and maintains the binding between agent and identifier through the agent's period of membership in the polity. We write $\Sig(A)$ for the identifier presently bound to agent $A$, understood as a commitment to a verification key together with the right to author signed messages under that key. The registry, not a fixed population, is the persistent backbone against which judgments are read; we develop this in \S\ref{ssec:registry} below.

\begin{definition}[Capability judgment]
A \emph{capability judgment} has the form
\[
\Gamma \vdash \Sig(A) : T
\]
where $T$ is a spatial-behavioral type and $\Gamma$ is a finite multiset of capability judgments $\Sig(B_i) : T_i$ representing the environmental dependencies on which $A$'s claim relies.
\end{definition}

The role of $\Gamma$ deserves emphasis. Where in conventional type theory $\Gamma$ records variable bindings local to a derivation, here $\Gamma$ records inter-agent dependencies. If Alice's ability to provide a database service depends on Bob's ability to provide reliable storage and Carol's ability to provide authentication, then Alice's judgment
\[
\Sig(\Bob):T_{\text{store}},\ \Sig(\mathsf{Carol}):T_{\text{auth}} \vdash \Sig(\Alice) : T_{\text{db}}
\]
becomes provable only by chaining in Bob's and Carol's own published judgments. Proof composition is the act of discharging environmental dependencies against the standing claims of other agents.

\subsection{The identity registry}\label{ssec:registry}

The treatment so far has spoken of identities as if they were fixed objects. In any realistic deployment this is too strong. Agents come and go. People are born and die. Devices are manufactured, deployed, retired. Software is updated beyond recognition. The population over which judgments range is dynamic, not fixed.

What \emph{is} fixed --- what must be fixed for the framework to denote anything stable across time --- is the system that maintains the binding of identifiers to agents. We call this an \emph{identity registry}: a persistent process responsible for issuing cryptographic identifiers to entering agents, maintaining the binding while the agent participates, and retiring or recycling the identifier when the agent departs.

\subsubsection*{Two familiar analogies}

An airline's baggage-claim system issues claim tags as travelers check bags, binds each tag to a traveler for the duration of a trip, and discharges the tag when the bag is reclaimed. The system persists across many travelers and many trips; no particular tag persists, but the institution that issues and reads them does. Similarly, a polity's voter-ID system issues voter identifiers to qualifying citizens, maintains them on the rolls, and removes them when eligibility lapses. The polity at any election is the set of identifiers currently active; the system that maintains the rolls is what makes the polity intelligible from one election to the next.

The identity registry of a reputation polity plays the same role. It is the persistent infrastructure underwriting the meaningful identification of agents over time, even though the agents themselves are not continuous.

\subsubsection*{The mobile-process-calculus perspective: unforgeability and the shard}

The institutional analogies above are useful but they leave the technical content of the registry abstract. A sharper structural identification is available in the mobile-process-calculus setting where the framework naturally lives.

In a mobile process calculus, the cryptographic identifier, the channel, and the object capability are not three different things in three different layers; they are the same primitive under three vocabularies. The signature key \emph{is} the capability. Holding the private key is, structurally, holding the unforgeable right to send as that identifier. The verification key is the public face of the channel on which messages from the holder are received as authentic. The framework's $\Sig(A)$ is therefore not a cryptographic apparatus bolted on top of the calculus; it is the calculus's existing name primitive with its cryptographic provenance made explicit.

The $\pi$-calculus and the rho-calculus differ sharply in how this primitive is treated.

The $\pi$-calculus's $\nu$-operator stipulates, axiomatically, that bound names are fresh, unique, and unforgeable from outside their scope: in $\nu x.P$, the name $x$ is asserted to be derivable from no construction the calculus's primitives can describe. Any \emph{deployment} of $\pi$ must produce names that are actually unforgeable in operational reality, but the calculus itself is silent about how. Public-key infrastructure, distributed consensus, secure enclaves --- whatever apparatus realizes the freshness claim --- is taken as given.

The rho-calculus is structurally honest about what $\pi$ conceals. There is no $\nu$-operator. Names are reified processes: $@P$ for some process $P$, structurally derived from what they name. Two names are distinct iff the underlying processes are structurally distinct. There is no primitive operation in rho that produces a name nobody else can construct; anyone who can reconstruct $P$ can derive $@P$. What rho's name primitive provides intrinsically is \emph{unguessability} rather than unforgeability --- names are computationally hard to find without knowing the construction, in the same sense that hash preimages are hard to find without knowing the input. The cryptographic analog is exact: rho's name discipline is, structurally, a commitment to one-way functions on processes.

The identity registry is precisely the apparatus that converts unguessability into operational unforgeability. By maintaining, under cryptographic assumptions plus consensus safety, the binding of unguessable names to agents and the record of which bindings are currently active, the registry yields operational unforgeability that is, under standard assumptions, indistinguishable in practice from $\pi$'s axiomatic version. The work is the same in both cases; what differs is whether the apparatus is foregrounded as an exhibited component (rho-with-registry) or buried in the calculus's stipulations ($\pi$-with-implicit-implementation).

\begin{remark}[The honest decomposition]
The rho-with-registry architecture is the honest decomposition of what $\pi$ packages opaquely. The institutional analogies --- baggage-claim, voter-ID --- name what the registry does institutionally; the technical analog under the metaphors is the conversion of unguessable names into operationally unforgeable ones.
\end{remark}

The composition is also clean with respect to the zero-knowledge refinements developed in \S\ref{sec:zk}: the registry provides operational unforgeability of names, and ZK adds operational hiding of relations between them. The two are independent layers of cryptographic discipline over the same underlying name structure. There is a Curry--Howard reading worth noting in passing: $\pi$'s $\nu x.P$ corresponds, propositionally, to an existential introduction whose witness is hidden --- ``there exists a name fresh in this scope.'' The rho-with-registry architecture realizes this existential witness computationally; the registry publishes the existential while the witness (the key, the construction of the underlying process) remains private to its holder. The framework's identity judgments inherit this discipline naturally, since they already place opaque witnesses (agent code) under typed claims; the registry simply extends the same discipline to the witness underlying the identifier itself.

\subsubsection*{Lifecycle operations}

A registry supports at least the following operations:

\begin{itemize}[leftmargin=*]
\item \textbf{Issuance.} A new agent acquires a freshly minted cryptographic identifier; the registry records the binding; the agent enters the polity and may begin signing judgments and attestations.
\item \textbf{Retirement.} A departing agent's identifier is removed from active status. Existing judgments and attestations bearing the identifier remain valid as historical artifacts attributable to the holder during the active period, but the identifier may no longer be used to sign new statements.
\item \textbf{Recycling (optional).} A retired identifier is reissued to a new agent. This is permissible only where the registry preserves temporal disambiguation: every signed statement bearing the identifier is attributable, by reference to the registry's records, to whichever holder was bound to the identifier at the time of signing.
\item \textbf{Revocation.} An identifier is forcibly removed from active status, typically in response to compromise of the underlying key. Revocation is sharper than retirement: it disclaims statements made after the compromise event, while leaving statements made before it valid.
\end{itemize}

These are operations of the registry, not of the reputation calculus proper. The reputation calculus consumes the registry's state as a precondition for the validity of signatures, and is otherwise agnostic to how the registry is implemented.

\subsubsection*{Consequences for the framework}

The registry-relative reading of identifiers ripples through the framework in a few specific ways. Each identifier $\Sig(A)$ should be understood as carrying an implicit time index --- the binding active at the moment a judgment or attestation is signed --- which we elide where context permits and restore where the lifecycle is at stake.

\begin{itemize}[leftmargin=*]
\item \emph{Provability is registry-relative.} A community-provable judgment in the sense of \S\ref{ssec:proofsys} is provable against a registry state. Subsequent changes to the registry may render previously valid derivations historically true but no longer actionable, in the same sense that a contract signed by a now-deceased party is historically valid but not the basis for new obligations.
\item \emph{Trace assembly is timestamped.} Each attestation in \S\ref{ssec:traces} is timestamped against the registry's notion of when the identifier-binding was active. This is what allows attestations referencing a since-recycled identifier to remain correctly attributable.
\item \emph{Reputation accrues to the binding, not to the bit-string.} An agent's reputation is a function of the judgments and attestations associated with bindings she has held. If a registry supports identifier rotation or migration (e.g., key compromise followed by re-issuance to the same agent under a new key), the reputation calculus follows the agent.
\item \emph{The registry is itself adjudicable.} Membership in the polity is governed by the registry. The registry's operations --- whom it admits, whom it retires, on what evidence --- are themselves subject to community adjudication via the same typed-claim apparatus. The governance of the polity reduces, in this sense, to the governance of its registry.
\end{itemize}

\begin{remark}[Identity registry as polity backbone]
The community of agents is dynamic; the registry is persistent. A reputation polity is defined not by a fixed population but by the persistence of the system that binds agents to identifiers and maintains those bindings through the agents' lifecycles. The framework's stability across time is the registry's stability across time.
\end{remark}

We leave the registry's implementation deliberately abstract. The framework is compatible with a range of substrates --- conventional public-key infrastructure, distributed ledgers, sharded namespaces, hierarchical federations of any of the above. What matters for the reputation calculus is that some such system exists, is persistent, and exposes the lifecycle operations above.

\subsection{The distributed proof system}\label{ssec:proofsys}

We assume a fixed sequent calculus over spatial-behavioral types --- determined by the OSLF construction applied to the underlying GSLT --- with the usual structural rules together with type-specific introduction and elimination rules. We augment this with two identity-specific rules.

\begin{prooftree}
\AxiomC{$\Sig(A)$ verifies on $\langle \text{``} \Gamma \vdash \Sig(A) : T \text{''} \rangle$}
\RightLabel{\textsc{Claim}}
\UnaryInfC{$\Gamma \vdash \Sig(A) : T$}
\end{prooftree}

\begin{prooftree}
\AxiomC{$\Gamma_1 \vdash \Sig(A) : T_1$}
\AxiomC{$\Gamma_2, \Sig(A) : T_1 \vdash \Sig(B) : T_2$}
\RightLabel{\textsc{Cut}$_{\text{id}}$}
\BinaryInfC{$\Gamma_1, \Gamma_2 \vdash \Sig(B) : T_2$}
\end{prooftree}

\textsc{Claim} embeds an agent's signed assertion as a leaf of the proof tree. The signature is what \emph{licenses} treating the assertion as a hypothesis available to the network. \textsc{Cut}$_{\text{id}}$ is the standard cut rule restricted to identity-indexed judgments; it discharges Alice's claim against the appearance of $\Sig(A) : T_1$ in another agent's environmental context. Composition of such cuts across the population is what assembles distributed proofs.

\begin{definition}[Provability]
The judgment $\vdash \Sig(A) : T$ is \emph{community-provable} if there exists a finite derivation in the augmented sequent calculus whose every leaf is a verified \textsc{Claim}. The set of such proofs is denoted $\Proofs(A,T)$.
\end{definition}

\begin{remark}[Structural choice: unrestricted hypotheses]\label{rem:unrestricted}
The calculus presented above is unrestricted in the sense of standard intuitionistic sequent calculi: hypotheses in $\Gamma$ are implicitly subject to weakening and contraction, and may be used any number of times in a derivation. This is the sound structural choice for reputation, because reputation is \emph{non-rivalrous}: Bob's published claim $\Sig(\Bob) : T_{\text{store}}$ is not depleted by being referenced in Alice's derivation, and may legitimately ground claims by many other agents simultaneously. The contraction rule reflects this fact about the underlying domain --- a hypothesis used twice is no different in standing from a hypothesis used once.

The extension of this framework to typed value transfer, developed in a companion position paper, requires a sharply different structural commitment. Value is rivalrous: a payment witness consumed in one transaction cannot be consumed in another, and any calculus that admitted contraction on value hypotheses would license double-spending. The natural setting is linear logic, with $\otimes$, $\multimap$, $\oplus$, and $\&$ tracking the resource bookkeeping that a currency requires. The full system is then a dual-context (DILL-style) calculus with an \emph{unrestricted zone} for capability claims about identities --- the calculus of the present paper --- and a \emph{linear zone} for payment witnesses, with weakening and contraction available on the first and disallowed on the second. The reputation framework presented here is, in this sense, the unrestricted half of a mixed linear/non-linear system whose other half handles value.
\end{remark}

\subsection{The trace semantics}\label{ssec:traces}

In parallel with the proof system, the community maintains a labelled transition system assembled from \emph{attestations} of communication events. An attestation is a signed statement of the form ``agent $B$ witnessed (or participated in) the event $e$ at time $t$'', where $e$ is a communication event in the rho-calculus sense: a synchronisation on a named channel carrying a payload. The attestation language and its discipline are developed in Section~\ref{sec:attest}.

Let $\mathcal{E}$ be the set of attested events, ordered by attestation timestamp into traces. We write $\sigma : A$ for a trace involving identity $A$ as participant.

\begin{definition}[Communal LTS]
The \emph{communal LTS} $\sem{\mathcal{E}}$ has agent identities as states and transitions $A \xrightarrow{e} A'$ wherever $A$ participates in $e$. The LTS is \emph{partial} (only attested edges appear), \emph{monotone} under attestation (new attestations only add edges), and \emph{multi-source} (edges may carry differing evidential strengths).
\end{definition}

OSLF-adequacy applied to the underlying GSLT yields, for each type $T$, a characteristic Hennessy--Milner formula $\phi_T$ such that a state inhabits $T$ (up to bisimulation) iff it satisfies $\phi_T$.

\begin{definition}[Refutation]
A trace $\sigma : A$ \emph{refutes} the claim $\Sig(A) : T$ if $\sigma$ witnesses a sub-state of $A$ satisfying $\neg \phi_T$, in the evidence-indexed HM logic of Section~\ref{ssec:modalities}.
\end{definition}

\subsection{Reputation as lattice position}\label{ssec:lattice}

The two processes --- distributed proof construction and communal trace assembly --- pull in opposite directions. Proof construction tries to certify; trace assembly tries to refute. Reputation is the joint state.

\begin{definition}[Reputation]
The \emph{reputation} of identity $A$ with respect to type $T$ is the pair
\[
\Reps(A,T) = \langle \Proofs(A,T),\ \Traces^{\bot}(A,T) \rangle
\]
where $\Traces^{\bot}(A,T)$ is the multiset of attested traces refuting the claim, each tagged with its modal-strength profile (Section~\ref{ssec:modalities}).
\end{definition}

The graded ordering is then immediate:

\begin{enumerate}[label=(\roman*)]
\item \emph{Proven}: $\Proofs(A,T) \neq \emptyset$ and $\Traces^{\bot}(A,T) = \emptyset$. The community has constructed a derivation; no attested trace refutes it.
\item \emph{Consistent}: $\Proofs(A,T) = \emptyset$ but $\Traces^{\bot}(A,T) = \emptyset$. No proof, but no refutation either. This is the strongest model-theoretic statement available short of provability.
\item \emph{Refuted to degree $\delta$}: $\Traces^{\bot}(A,T) \neq \emptyset$, with $\delta$ a measure built from the multiset and its modal-strength profile.
\item \emph{Provably false}: a derivation exists for $\vdash \Sig(A) : \neg T$, or for an incompatible $T'$.
\end{enumerate}

\begin{remark}[Anti-scalarization]
We emphasise --- following the principle that \emph{typeless scalars offer no objective adjudication process, because the process the scalar witnesses is unrecoverable from the scalar itself} --- that $\Reps(A,T)$ is a \emph{structured} object. Different communities, with different decision problems, project this structure to scalars (or richer decision-relevant objects) by different aggregation functions. The aggregation is exposed as a community-specific commitment, adjudicable in its own right, not buried inside an opaque score.
\end{remark}

\begin{remark}[Coverage-indexing]
The grading is necessarily indexed by community observational coverage. An agent in a sparsely observed corner of the network earns ``consistency'' cheaply. We discuss attestation economics --- the phlogiston-metered ``valuable friction'' that makes coverage non-trivial --- in Section~\ref{ssec:friction}.
\end{remark}

\section{Attestations and the rho capability discipline}\label{sec:attest}

\subsection{Three tiers}

Attestations stratify into three tiers, each defined by the attestor's epistemic relationship to the event:

\begin{definition}[Attestation tiers]
For an event $e$ on channel $c$ carrying payload $p$ between principals $A$ and $A'$:
\begin{itemize}[leftmargin=*]
\item $B \AttD e$ (\emph{direct}): $B \in \{A, A'\}$; $B$ executed the send or receive on $c$.
\item $B \AttW e$ (\emph{witnessed}): $B \notin \{A,A'\}$ but $B$ holds the name $c$ and observed the synchronisation.
\item $B \AttH^{\chi} e$ (\emph{hearsay}): $B$ does not hold $c$ but was told of $e$ by some other attestor, with $\chi$ recording the full provenance chain of intermediaries.
\end{itemize}
\end{definition}

The hearsay tier requires the full provenance chain rather than a flat \emph{hearsay} tag. A bare hearsay flag is forgeable through cycles: $E$ tells $B$, $B$ tells $D$, $D$ tells $E$, and there now appear to be three independent hearsay attestations of a fact with a single origin. Explicit chains make the dependency structure visible and let the reputation calculus discount accordingly.

\subsection{Capability alignment}

The tier discipline is not external constraint imposed on the rho calculus. It \emph{is} the rho calculus's name-as-capability discipline applied at the meta-level.

\begin{proposition}[Capability alignment]
In a rho-calculus setting, the right to attest at tier $\AttD$ for an event on channel $c$ is precisely the right to send or receive on $c$. The right to attest at tier $\AttW$ for an event on $c$ is precisely possession of $c$ as a name. The right to attest at $\AttH$ requires no nominal access, only a verified hearsay chain.
\end{proposition}

Consequently the \emph{witnessing radius} of any event is exactly the distribution of its channel name. The social topology of who-can-attest-to-whom is the nominal topology of who-shares-names-with-whom. Reputation observability and capability locality are the same structure under two interpretations.

\subsection{Evidence-indexed modalities}\label{ssec:modalities}

The OSLF functor, applied to the underlying GSLT, generates a Hennessy--Milner-style logic with diamond modalities $\langle a \rangle \phi$. Stratifying attestation refines these into an \emph{evidence-indexed} family:

\begin{align*}
\langle a \rangle^{\mathsf{D},B}\phi &\quad\text{``some attestor $B$ directly attests an $a$-transition into a state satisfying $\phi$''} \\
\langle a \rangle^{\mathsf{W},B}\phi &\quad\text{``some attestor $B$ witnessed an $a$-transition into a state satisfying $\phi$''} \\
\langle a \rangle^{\mathsf{H},\chi}\phi &\quad\text{``the hearsay chain $\chi$ supports an $a$-transition into a state satisfying $\phi$''}
\end{align*}

Modal strength is partially ordered:
\[
\modD{a} \;\succeq\; \modW{a} \;\succeq\; \langle a \rangle^{\mathsf{H},\chi}
\]
with the hearsay tier further refined by chain length and signing density along $\chi$.

\begin{definition}[Refutation strength]
The strength of a refuting trace $\sigma$ for the claim $\Sig(A):T$ is the meet, in the modal-strength order, of the modalities appearing in the satisfying assignment of $\neg \phi_T$ at $\sigma$. A refutation is only as strong as its weakest evidential link.
\end{definition}

This gives a principled, type-driven grading of refutations, and is exactly the structure that resists collapse into a scalar. The full refutation profile of an identity is a partial-order-valued object; projection to a scalar is a community-specific choice with its own adjudication.

\subsection{Signing-on-send as valuable friction}\label{ssec:friction}

A unilateral attestation $B \AttD e$ where $A$ is the counterparty raises the obvious adversarial question: what stops $B$ from lying about $A$? The framework's answer is structural rather than juridical.

\begin{proposition}[Non-repudiation by signed payload]
If the payload $p$ of event $e$ carries $A$'s signature, then $B$'s attestation $B \AttD e$ extends to a non-repudiable record: $B$ produces $A$-signed evidence that $A$ engaged in the event.
\end{proposition}

Where $A$ has not signed, $B$'s direct attestation degrades to a one-sided claim, which $A$ may contradict with her own attestation of the period in question, and the reputation calculus grades both as conflicting evidence rather than treating either as decisive.

The implication is that signing-on-send constitutes \emph{valuable friction at the attestation layer}: cheap unsigned communications generate cheap, weak evidence; expensive signed communications generate strong, non-repudiable evidence. Agents self-select their evidential exposure by choosing how much to sign --- which is the same currency in which they self-select their economic exposure on the rho calculus's phlogiston metering of computation.

\section{Difference from ordinary reputation systems}\label{sec:diff}

We highlight the structural differences from the typeless-scalar reputation systems prevalent on existing platforms.

\paragraph{Adjudicability.} A scalar rating ($4.3 \pm 0.2$ stars; karma; ELO) is the output of an aggregation process not exhibited at the rating site. Disputes over the rating cannot be adjudicated because the process the scalar witnesses is unrecoverable from the scalar. In the present framework, every reputation claim decomposes into structurally inspectable components: a set of signed proofs, a set of attested traces with full provenance, and a community-specific projection. Each is adjudicable in its own right.

\paragraph{Type-conformance vs.\ aggregated behavior.} Scalar systems aggregate heterogeneous behavior into a single dimension. The present framework asks a more specific question: does $\Sig(A)$ inhabit type $T$? This is meaningful only with respect to a type that describes the relevant service. Reputation thereby tracks \emph{capability}, not popularity or compliance.

\paragraph{Sybil resistance.} The use of cryptographic identity at the type-judgment level means a Sybil army of fresh identities has no reputation, not merely a low one. Reputation accrues only through accumulation of signed evidence over time, and the cost of generating that evidence is bounded below by the cost of the corresponding signed communications.

\paragraph{Manipulation surface.} The manipulation surface of a scalar system is the aggregation function plus its input channel: stuff the ballot, brigade the reviews, game the algorithm. The manipulation surface of the present framework is forging signatures (cryptographically bounded), constructing fraudulent proofs (rejected by the type system), or generating fraudulent traces (countered by the capability alignment of attestation rights). None of these admits cheap mass attack.

\paragraph{Compositionality.} Scalar ratings compose poorly: there is no principled way to combine ``$A$ is rated $4.3$ on platform $X$'' with ``$A$ is rated $\beta$ on platform $Y$''. The judgment $\Gamma \vdash \Sig(A) : T$ composes by cut. Cross-community reputation transfer is a matter of importing signed judgments and re-verifying their dependencies, which is structurally clean.

\paragraph{Privacy.} Scalar systems typically reveal the rating but conceal the basis. The present framework conceals the agent's internal code (it was never exposed) while exposing the basis. Section~\ref{sec:zk} discusses how zero-knowledge techniques recover privacy of the basis where this is desired.

\section{MeTTaIL implementation sketch}\label{sec:mettail}

MeTTaIL (rholang~1.4) is engineered precisely for this kind of construction: it provides the means to define a language as a triple of grammar (syntactic categories), equations on types, and rewrites on types, and supports sending and receiving terms of the defined language across the network. We sketch the reputation system as such a defined language.

\subsection{Grammar}

The reputation language's grammar is given in BNFC-like syntax:

\begin{lstlisting}[language=rholang,caption={Reputation language grammar (BNFC-style)}]
-- Identities and signatures
Identity. Id ::= "Sig" "(" AgentName ")" ;

-- Spatial-behavioral types (parametric on the underlying GSLT)
TypeAtom. Type ::= TypeAtom ;
TypeConj. Type ::= Type "&" Type ;
TypePref. Type ::= Action "." Type ;
TypeStar. Type ::= "!" Type ;
-- ... (full GSLT type grammar elided)

-- Capability judgments
EmptyCtx. Ctx ::= ;
ConsCtx.  Ctx ::= Ctx "," Id ":" Type ;
Judgment. Jud ::= Ctx "|-" Id ":" Type ;

-- Signed claims (Section 2.2)
Claim. SignedClaim ::= "claim" "(" Id "," Jud "," Sig ")" ;

-- Attestations (Section 3.1)
EvtSend. Event ::= "send" "(" ChanName "," Payload ")" ;
EvtRecv. Event ::= "recv" "(" ChanName "," Payload ")" ;
EvtSync. Event ::= Id "<->" Id "@" ChanName "(" Payload ")" ;

AttD. Attestation ::= "attD" "(" Id "," Event "," Time "," Sig ")" ;
AttW. Attestation ::= "attW" "(" Id "," Event "," Time "," Sig ")" ;
AttH. Attestation ::= "attH" "(" Id "," Event "," Chain "," Time "," Sig ")" ;

-- Hearsay provenance chains
ChainNil.  Chain ::= "[]" ;
ChainCons. Chain ::= Attestation "::" Chain ;

-- Refutation profiles and reputation objects
ModD. Strength ::= "D" "(" Id ")" ;
ModW. Strength ::= "W" "(" Id ")" ;
ModH. Strength ::= "H" "(" Chain ")" ;

Trace.   Trace      ::= "[" [Attestation] "]" ;
RepObj.  Reputation ::= "rep" "(" Id "," Type "," [SignedClaim] "," [Trace] ")" ;
\end{lstlisting}

\subsection{Equations}

Equations identify terms that should be regarded as denoting the same object under the type-theoretic and cryptographic disciplines:

\begin{lstlisting}[language=rholang,caption={Equations on the grammar}]
-- Signature extensionality: signatures verify-equal iff they verify the same payload.
sig_ext :  Sig(A, m1) == Sig(A, m2)   when   payload(m1) == payload(m2) ;

-- Context permutation: contexts are multisets.
ctx_perm:  (Gamma1, j1, j2)  ==  (Gamma1, j2, j1) ;

-- Hearsay chain de-duplication: collapse cycles.
chain_dedup:  c1 ++ [a] ++ c2 ++ [a] ++ c3  ==  c1 ++ [a] ++ c2 ++ c3 ;

-- Direct dominates witnessed for the same event.
att_dominate: attW(B,e,t,s)  ==  attD(B,e,t,s)   when   B in participants(e) ;
\end{lstlisting}

\subsection{Rewrites}

Rewrites give the operational behaviour of the proof system and trace assembly:

\begin{lstlisting}[language=rholang,caption={Proof-construction and trace-extension rewrites}]
-- Claim: a verified signed claim becomes a leaf.
claim_intro:
    claim(Sig(A), (Gamma |- Sig(A) : T), s)
    -->
    Gamma |- Sig(A) : T
    when verify(s, A, "Gamma |- Sig(A) : T") ;

-- Cut_id: discharge an environmental dependency.
cut_id:
    (Gamma1 |- Sig(A) : T1),  (Gamma2, Sig(A) : T1 |- Sig(B) : T2)
    -->
    (Gamma1, Gamma2 |- Sig(B) : T2) ;

-- Trace extension: a new attestation extends the communal LTS.
trace_ext:
    Trace([as]),  Attestation(a)
    -->
    Trace([as, a])
    when consistent_time(a, as) ;

-- Refutation detection (delegated to HM-model-checker).
refute:
    rep(A, T, Cs, Ts),  Trace(sigma)
    -->
    rep(A, T, Cs, Ts ++ [sigma])
    when sigma  |=  neg(charformula(T))  in  modal_logic(T) ;
\end{lstlisting}

\subsection{Hosting in rholang processes}

The defined language is then operationalised by ordinary rholang processes which carry typed terms on well-known channels. A schematic outline:

\begin{lstlisting}[language=rholang,caption={Rholang processes hosting the reputation language}]
new claimsCh, attestCh, repQueryCh, repAnswerCh in {
  // Claim ingress: receives signed claims, validates, deposits in claim store.
  contract claimsCh(@msg) = {
    match msg with
      claim(sig, jud, s) =>
        if verify(s, agentOf(sig), jud) then storeClaim(claim(sig, jud, s))
        else Nil
  } |

  // Attestation ingress: similar discipline for attestations.
  contract attestCh(@msg) = {
    match msg with
      attD(B, e, t, s) | attW(B, e, t, s) | attH(B, e, chi, t, s) =>
        if verify(s, B, payloadOf(msg)) then storeAttestation(msg)
        else Nil
  } |

  // Reputation queries: a community member asks for rep(A, T).
  contract repQueryCh(@A, @T, @replyCh) = {
    new proofs, traces in {
      gatherClaims(A, T, proofs) |
      gatherTraces(A, T, traces) |
      for(@p <- proofs; @ts <- traces) {
        replyCh!(rep(A, T, p, ts))
      }
    }
  }
}
\end{lstlisting}

The HM model-checker invoked by the \texttt{refute} rewrite is itself a rholang contract, automatically generated by the OSLF functor from the GSLT defining the spatial-behavioral types. This is the structural point of using MeTTaIL: the language definition, the model-checker, the proof system, and the attestation infrastructure are all uniformly rho processes operating on uniformly defined terms.

\subsection{Phlogiston economics}

Each attestation, claim, and proof step incurs phlogiston cost, paid by the publishing agent. This is the framework's coverage-resistance mechanism. Sparse coverage of an agent's behaviour is structurally cheap to remedy --- anyone holding the relevant channel names can attest --- but \emph{costly} in phlogiston. Communities concerned about an agent in their corner of the network can subsidise observation; the cost is paid in the same metering used for ordinary computation.

\section{Connection to zero-knowledge proofs}\label{sec:zk}

The framework's central conceit --- adjudicating capability without inspecting the agent's code --- has a strong structural parallel with zero-knowledge proof systems, where a verifier becomes convinced that a prover knows a witness $w$ satisfying a relation $R(x,w)$ without learning $w$. The parallel is not metaphorical: ZK techniques compose with the framework at several levels.

\subsection{Identity-as-witness as analogue of statement-witness}

In a ZK proof system, $(x,w) \in R$ is the relation, $x$ is public, $w$ is hidden. In the reputation framework, $(\Sig(A), C_A) \models T$ is the relation (Alice's internal code $C_A$ inhabits the type $T$), $\Sig(A)$ and $T$ are public, and $C_A$ is hidden. The community plays the role of verifier and is convinced by accumulating signed evidence and the absence of refuting traces.

The disanalogy is interesting: ZK gives \emph{certainty} of witness existence by proof-theoretic means; the reputation framework gives \emph{graded confidence} of type conformance by a hybrid proof-theoretic and observational means. ZK is sharper but narrower. The framework's grading is exactly the price of admitting observational evidence into the adjudication.

\subsection{Attestations as ZK statements}

The framework's attestations need not reveal everything they currently reveal. Several layers of ZK refinement are natural:

\paragraph{Channel-name-hiding witnessing.} A witnessing attestation $B \AttW e$ currently reveals the channel $c$ on which $e$ occurred. But channel names are capabilities; revealing them is leaking authority. A more disciplined attestation publishes a ZK proof that ``$B$ holds some channel $c$ such that $e$ occurred on $c$'', without revealing $c$. This preserves the witnessing semantics --- $B$ has demonstrated capability-justified observation --- while maintaining nominal locality.

\paragraph{Payload-content-hiding attestation.} An attestation of an event with sensitive payload can publish ZK evidence that the payload satisfies some predicate (e.g., that it is well-formed under the relevant type) without revealing the payload itself. This is the right discipline for, e.g., medical-record-bearing attestations in the VA CRADA context, where capability conformance is to be adjudicated without exposing patient data.

\paragraph{Selective-disclosure provenance.} Hearsay chains $\chi$ currently expose the full chain of intermediaries. A ZK refinement publishes a proof that ``a valid chain of length $\leq n$ exists between origin and current attestor'', without exposing the intermediate identities. This is appropriate where intermediaries' identities are themselves sensitive.

\subsection{Type-conformance proofs as ZK certificates}

A community-provable judgment $\vdash \Sig(A) : T$ has a witnessing proof tree composed from signed claims of multiple agents. The proof tree may itself be sensitive --- it exposes Alice's environmental dependency graph, which is competitive intelligence. A ZK refinement publishes a proof that ``a valid proof tree exists in the augmented sequent calculus'', without exposing the tree.

The compositional analogue is incrementally verifiable computation (IVC) and recursive SNARK constructions: as new claims are published and the proof tree grows, the ZK certificate is incrementally updated rather than recomputed. Reputation snapshots become succinct, verifiable, and privacy-preserving certificates of capability over time.

\subsection{Curry--Howard angle}

Under the Curry--Howard correspondence, proofs are programs; under the ZK refinement, the programs' constructive content is concealed but their type is verified. This is precisely the discipline already at work in the framework's identity judgments: $\Sig(A)$ is a term whose constructive content (the code) is concealed, but the type is being verified by the network. The ZK refinement of proof certificates simply extends this discipline from the leaf level (where $\Sig(A)$ already conceals $C_A$) to the entire proof structure.

\subsection{What ZK does not provide}

It is worth noting clearly what ZK does \emph{not} replace. ZK proofs adjudicate witness existence under a fixed relation. They do not adjudicate \emph{contested observation}, which is what the trace-refutation side of the framework handles. If Bob attests in ZK that $e$ occurred and Alice attests in ZK that her behaviour at $t$ was incompatible with $e$, the community still requires a non-ZK adjudication of the conflict, which the lattice-structured reputation object supports. ZK refines the privacy of evidence; it does not subsume the adjudicative function of communal observation.

\section{Conclusion and further directions}

The framework rests on a single move: putting an opaque cryptographic identity in the term position of a typing judgment, and using the GSLT/OSLF apparatus to adjudicate the resulting claim via two coupled processes. The yield is a reputation object that is structurally adjudicable, type-driven rather than scalar, naturally Sybil-resistant, compositional under cut, and which composes cleanly with zero-knowledge refinements.

Several directions invite continuation. The aggregation question --- how a specific community projects the structured reputation object onto a decision --- connects directly to the Coalition Logic and consensus-synthesis line of work on semitopologies, with proof terms corresponding to join patterns. The phlogiston economics of attestation invites detailed treatment as an extension of the existing rholang metering. The VA CRADA application provides a concrete setting in which the ZK refinements of Section~\ref{sec:zk} are not optional but required by the regulatory context. The governance of the identity registry --- whom the polity admits, on what evidence, and through what process of communal adjudication --- is a separate axis of investigation that the framework opens but does not settle in the present paper. And the broader question --- whether observer-generated logics over GSLTs furnish a foundation for reputation across the full F1R3FLY property portfolio (F1R3Sky, F1R3Eats, F1R3Tunes, F1R3Games) --- becomes tractable once reputation is type-indexed and community-projected rather than platform-scalar.

\bigskip

\noindent\textit{Acknowledgments.} This document was developed in dialogue. The substantive framework, including the inversion of the term position, the proof-system/trace duality, the three-tier attestation hierarchy, the anti-scalarization principle, and the MeTTaIL hosting target, is Meredith's. Claude contributed structural articulation, the evidence-indexed modal refinement, and the elaboration of the ZK parallels.

\end{document}
